\subsection{\problemblock{*Aero} Data Block}\label{sec:block-aero}

The \problemblock{*aero} data block is used to define the aerodynamic
coefficients. It is optional; if it is not input, all aerodynamic coefficients are set
to zero. If more than one \problemblock{*aero} data block is input in a segment,
the aerodynamic forces are computed for each \problemblock{*aero} data block and
then added together.

There are three different sets of aerodynamic coefficients that can be input: the
axial and normal force coefficients ($C_A$ and $C_N$), the lift, drag, and side-force
coefficients ($C_L$, $C_D$, and $C_S$), and the body $x$-, $y$-, and $z$-axis force
coefficients ($C_X$, $C_Y$, and $C_Z$). The aerodynamic coefficients are defined
in \cref{sec:aerodynamic-forces}.

Within each \problemblock{*aero} data block, a consistent set of aerodynamic
coefficients must be input. Coefficients from different sets cannot be mixed. For
example, $C_L$ and $C_D$ are in the same set, so they form a valid set of
coefficients, whereas $C_L$ and $C_A$ are not in the same set, so they cannot be
used together.

If a segment has more than one \problemblock{*aero} block, it is recommended
that all of them use the same type of aerodynamic coefficients. If this is done, then
all of the output aerodynamic coefficients are the sum of the coefficients from each
\problemblock{*aero} block. Otherwise, the forces used to compute the trajectory
are correct, but only the $C_L$, $C_D$, and $C_S$ output coefficients are correctly
totaled.

\taossubhead{Constant Coefficients}

Aerodynamic coefficients can be input either as constants or as tables. Constant
values are input by giving the coefficient name and its value as follows:

\begin{lstlisting}[style=taosmanual]
*aero ca=0.0215 cn=0.118 sref=5.437
\end{lstlisting}

This data block defines the axisymmetric coefficients $C_A$ and $C_N$. These
values are constant for the duration of the segment.

Besides the aerodynamic coefficient values, a value has been given for the reference
area $S_{\mathrm{ref}}$. The reference area is used to compute the aerodynamic
forces from the coefficients and it has default units of \si{\square\foot}. It is
required when all coefficients are set to constants; when tables are used, the
reference area can be given in one of the tables.

When a reference area is specified in the \problemblock{*aero} data block, it
overrides values given in tables. When multiple \problemblock{*aero} blocks are
given or when aerodynamic tables contain different values of $S_{\mathrm{ref}}$,
\statevar{sref} should not be given in the \problemblock{*aero} data block.

\taossubhead{Aerodynamic Tables}

Aerodynamic coefficients can also be given in tables by setting the coefficient names
equal to a table identification name enclosed in parentheses. For example,

\begin{lstlisting}[style=taosmanual]
*aero ca=(ca-clean) cn=(cn-clean)
\end{lstlisting}

requests a table called \tableid{ca-clean} for the axial-force coefficient and a table
called \tableid{cn-clean} for the normal-force coefficient.

Tables can be used for some coefficients while others can be set to constants. In the
following example, the lift coefficient is set to a table name and the drag coefficient
is set to a constant:

\begin{lstlisting}[style=taosmanual]
*aero cl=(lift-x) cd=0.2351
\end{lstlisting}

The side-force coefficient is not given, so it is set to zero. The reference area is not
given either, so it must be provided in the \tableid{lift-x} table.

\taossubhead{User-Defined Variables}

Tables can be functions of user-defined variables as shown in
\cref{sec:table-user-defined-variables}. If user-defined variables are
referenced in a table, then values for them must be given in the problem file.

One place values can be given for user-defined variables is in the
\problemblock{*aero} data block. Although values can be given for any user-defined
variable within this data block, usually values are only given for the user-defined
variables in the aerodynamic tables. Values for other user-defined variables are
given in the data blocks where they are referenced. This way the table name and its
user-defined variables are given together in the problem file. Another method is to
give values for the user-defined variables in a \problemblock{*constants} data block
(\cref{sec:block-constants}).

Once a value has been given for a user-defined variable, it remains set at that value
unless changed in a subsequent data block. This means that values for user-defined
variables do not have to be repeated within each segment if they do not change.
However, it is good practice to provide values in each segment to avoid input errors.

The following \problemblock{*aero} data block defines values for two user-defined
variables:

\begin{lstlisting}[style=taosmanual]
*aero cd=(drg-model) cl=(lft-model)
  flaps=25, gear=0, sref=225.0
\end{lstlisting}

This block provides table names for $C_L$ and $C_D$, and it sets the reference area
to \SI{225}{\square\foot}. It also provides values for two user-defined variables:
\texttt{flaps} and \texttt{gear}. These variables must appear in the tables;
arbitrary variable names that never appear in a table cannot be input. These
variables could be flags in the table used for flap deflection and gear up or down.
